\documentclass[11pt,reqno]{amsart}

% ---------------------------------------------------------------------------
%  From minimal description to E6: a parameter-free chain through the
%  figure-eight knot complement.
%
%  Self-contained for arXiv: no \input of anything outside this directory, no
%  local paths, no .bib (bibliography is inline, see \thebibliography).
%  Gate 5: no physical identification is made in the body.
% ---------------------------------------------------------------------------

\usepackage{amsmath,amssymb,amsthm,mathtools}
\usepackage[T1]{fontenc}
\usepackage{lmodern}
\usepackage{microtype}
\usepackage{booktabs}
\usepackage{enumitem}
\usepackage[margin=1.1in]{geometry}
\usepackage{xcolor}          % REQUIRED before hyperref for the `!' colour mixing below
\usepackage[colorlinks=true,linkcolor=blue!60!black,citecolor=green!45!black,urlcolor=blue!70!black]{hyperref}

\theoremstyle{plain}
\newtheorem{theorem}{Theorem}[section]
\newtheorem{proposition}[theorem]{Proposition}
\newtheorem{lemma}[theorem]{Lemma}
\newtheorem{corollary}[theorem]{Corollary}
\theoremstyle{definition}
\newtheorem{census}[theorem]{Census}
\newtheorem{nogo}[theorem]{No-go}
\theoremstyle{remark}
\newtheorem{scope}[theorem]{Scope}
\newtheorem{remark}[theorem]{Remark}

\newcommand{\Q}{\mathbb{Q}}
\newcommand{\Z}{\mathbb{Z}}
\newcommand{\C}{\mathbb{C}}
\newcommand{\R}{\mathbb{R}}
\newcommand{\SL}{\mathrm{SL}}
\newcommand{\SU}{\mathrm{SU}}
\newcommand{\su}{\mathfrak{su}}
\newcommand{\so}{\mathfrak{so}}
\newcommand{\uu}{\mathfrak{u}}
\newcommand{\ee}{\mathfrak{e}}
\newcommand{\tr}{\operatorname{tr}}
\newcommand{\Vol}{\operatorname{Vol}}
\newcommand{\disc}{\operatorname{disc}}

\begin{document}

\title[From minimal description to $E_6$]{From minimal description to $E_6$:\\
a parameter-free chain through the figure-eight knot complement}

\author{Drit\"ero Mehmetaj}
\address{Independent Researcher}

\date{\today}

\subjclass[2020]{Primary 57K10, 11R16, 57M50; Secondary 17B25, 17A75, 11J06}
\keywords{figure-eight knot complement, Sturmian words, continued fractions,
Lagrange spectrum, McKay correspondence, exceptional Jordan algebra,
cone-manifold geometric transition, parameter-free derivation}

\begin{abstract}
We fix a single principle---that a description cannot terminate---and follow it
without importing anything from outside, recording what it costs. With three
declared choices, whose alternatives we compute and report, the principle yields a
specific hyperbolic $3$-manifold: the figure-eight knot complement. Its own
arithmetic then determines a $78$-dimensional exceptional Lie algebra with a
distinguished four-element charge torus, and a cascade of centralizers of that
torus terminates, exactly, at $\su(3)\oplus\su(2)\oplus\uu(1)^3$. No measured
quantity enters any computation.

We report the negative half with equal weight. The construction is provably silent
about magnitudes: the object's amphichirality forces its Chern--Simons invariant to
vanish, deleting the only term in its boundary action that could carry a scale.
Five predictions were sealed in advance and compared against measurement; all five
missed, and the surface of comparisons the construction licenses was subsequently
enumerated and found exhausted. We state the arithmetic results as mathematics and
make no physical identification.
\end{abstract}

\maketitle

\section{Introduction}\label{sec:intro}

Fix a principle: \emph{a description that cannot terminate}. Ask what the least such
description looks like, follow the answer without importing anything from outside,
and record what it costs.

This paper carries out that programme and reports both halves of the result. From
the principle, with three declared choices whose alternatives we compute, one
reaches a specific hyperbolic $3$-manifold---the figure-eight knot complement. Its
own arithmetic then determines a $78$-dimensional exceptional Lie algebra with a
distinguished four-element charge torus, and a cascade of centralizers of that
torus terminates, exactly, at $\su(3)\oplus\su(2)\oplus\uu(1)^3$ with the global
form $[\SU(3)\times\SU(2)\times \mathrm{U}(1)]/\Z_6$. \textbf{No measured quantity
enters any computation at any point.}

The second half we ask the reader to weigh equally. The same construction is
\emph{provably} silent about magnitudes (\S\ref{sec:negatives}), and five
pre-registered comparisons with measurement all missed.

\subsection{What is claimed}
A \emph{cost theorem}: three declared choices, no fitted parameters, and no measured
input yield the structural data of \S\S\ref{sec:frame}--\ref{sec:realform}. We do
\textbf{not} claim a derivation of physics. Internal names for mathematical objects
are defined at first use.

\subsection{What is unusual, stated plainly}
Of the $43$ links in the chain, $39$ are theorems, identities, censuses or no-go
results; four are declared choices. Three of the four occur \emph{before} the
manifold appears, and one occurs after the algebra is in hand. Between the manifold
and the algebra there is not one declared choice.

\subsection{What the reader should be suspicious of}\label{sec:suspicious}
Two things, addressed where they arise rather than in a discussion section.

First, reaching this particular exceptional algebra is \emph{generic}: roughly one
hyperbolic $3$-manifold in three, and five of seven grammars in the relevant
family, arrive at it. \textbf{Arrival there is not evidence.} What is not generic
is the \emph{entrance} (\S\ref{sec:entrance}). We can be quantitative rather than
rhetorical, because the question was settled by enumeration.

\begin{census}\label{cen:gamma}
Of the $20$ steps below the entrance, $0$ consume the manifold's fundamental group,
$18$ are standard algebra downstream of the McKay assignment, and exactly $1$
consumes the arithmetic---the entry itself, where the Kleinian trace field selects
\emph{which} McKay group. None is undecidable.
\end{census}

So the paper's structure mirrors a measured fact: everything after the entrance is
exceptional-algebra theory checkable against the standard references, and the single
step where this object's arithmetic does work is the one \S\ref{sec:entrance} is
about. A reader who wishes to attack the paper should attack that step; attacking
the cascade attacks Slansky~\cite{Slansky}.

Second, the cost claim is a claim about what entered the computations, and prose
cannot establish it. Appendix~\ref{app:verify} supplies one runnable verification
per numerical claim.

\section{The principle and its price}\label{sec:price}

We take as starting point that description is inexhaustible: no finite description
closes. The mathematical content is a minimality condition (\S\ref{sec:genesis});
the metaphysical content is a choice, and we price it. \textbf{A declared choice
whose alternatives are not computed is a free parameter wearing a different name},
so the computations are part of the claim.

\begin{description}[leftmargin=1.6em,style=nextline]
\item[(C3) Description is inexhaustible.] The one metaphysical commitment.
  \emph{Price.} Both alternatives degenerate. The periodic counterpart collapses:
  the whole $\det=+1$, $|\tr|\le 2$ family is finite-order or reducible---no
  pseudo-Anosov element, no hyperbolic carrier. Neither alternative supports the
  construction at all.
\item[(C4) The word is realized on the once-punctured torus.] A geometric carrier
  over combinatorial ones. \emph{Price, and it is consequential.} The canonical
  non-geometric carriers---the tiling hull, and the Effros--Shen algebra with
  $K_0=\Z[\varphi]$---see only the real quadratic data. All four pre-registered
  redundancy witnesses fail exactly: in particular $x^2+3$ remains irreducible over
  $\Q(\sqrt5)$. \textbf{The imaginary quadratic field $\Q(\sqrt{-3})$ is bought at
  geometrization and nowhere earlier.}
\item[(C5) The monodromy is orientation-preserving.] Equivalently, the golden slope
  is squared. \emph{Price, and it is the largest.} \textbf{The discarded alternative
  is not degenerate.} It is the Gieseking manifold---the orientation double cover's
  parent of the manifold we keep, with dilatation $\varphi$ upstairs against
  $\varphi^2$ below.
\end{description}

\begin{table}[h]
\centering
\begin{tabular}{@{}ll@{}}
\toprule
choice & status of the alternative\\
\midrule
C3 & degenerates---no carrier exists\\
C4 & survives, but sees only the real quadratic data\\
C5 & \textbf{survives intact as a named manifold}\\
\bottomrule
\end{tabular}
\end{table}

\textbf{C5 is the honest scar in the argument.} It is a fork between two genuine
objects decided by a stipulation, and no computation in this paper forces it. We
flag it here rather than in a limitations section because a reader who discovers it
late is entitled to distrust everything before it.

\section{From minimal description to a knot}\label{sec:genesis}

\begin{theorem}[Morse--Hedlund~\cite{MorseHedlund}]\label{thm:mh}
If an infinite word $w$ is aperiodic then its factor complexity satisfies
$p(n)\ge n+1$ for every $n\ge1$. Words attaining $p(n)=n+1$ for all $n$ are exactly
the Sturmian words.
\end{theorem}

Periodic words close, so under the principle they are excluded; among the words
that do not close, Theorem~\ref{thm:mh} identifies a unique minimal complexity
class, and it is non-empty.

Sturmian words are classified by a single irrational slope $\alpha$. The slope is
itself a parameter of the description, and applying the principle to \emph{it} is
what removes the last freedom. The relevant notion of ``least reducible to a
terminating description'' is worst rational approximability, measured by the
Lagrange number $L(\alpha)$.

\begin{theorem}[self-selection]\label{thm:self}
Among slopes with eventually constant continued fraction $[a;a,a,\dots]$, $a\ge1$,
one has $L=\sqrt{a^2+4}$. This is strictly increasing in $a$; the unique minimiser
is $a=1$, the golden slope $\varphi$, with $L(\varphi)=\sqrt5$, the bottom of the
Lagrange spectrum (Hurwitz).
\end{theorem}

Two features make this a theorem rather than a preference. \emph{(i) It is a fixed
point, not a selection}: a continued fraction all of whose partial quotients equal
$1$ is precisely a fixed point of $t\mapsto 1+1/t$, so the statements ``the
continued fraction is all-$1$s'' and ``$x=1+1/x$'' are the same statement.
\emph{(ii) The minimiser is unique}, by strict monotonicity.

\begin{scope}\label{sc:markov}
\textbf{Nothing in this section is offered as new.} The extremality is classical
(Hurwitz; Markov), and the whole setting---Sturmian and Christoffel words, continued
fractions, and Markoff theory, together with Sturmian morphisms and the positive
automorphisms of $F_2$---is the subject of Reutenauer's monograph~\cite{Reutenauer},
the link back to Frobenius (1913). What is ours is the \emph{use} made of them:
applying the minimality principle to its own parameter, which is a modelling step,
not a theorem.
\end{scope}

By C4 the word is realized on the once-punctured torus, and by C5 the monodromy is
$M=\left(\begin{smallmatrix}2&1\\1&1\end{smallmatrix}\right)$, $\tr M=3$.

\begin{theorem}[Thurston; Riley]\label{thm:thurston}
The mapping torus of the once-punctured torus under $M$ is the complement of the
figure-eight knot. It carries a unique complete hyperbolic structure, and its
Kleinian trace field is $\Q(\sqrt{-3})$.
\end{theorem}

This ends the genesis. \textbf{Three declarations have been spent; two theorems and
one classical realization carried the rest; no numerical value has entered.}

\section{The object's own arithmetic}\label{sec:arith}

Everything here is a property of the manifold of \S\ref{sec:genesis}; nothing is
imported and no choice is made.

\begin{theorem}\label{thm:v4}
The manifold's intrinsic arithmetic forces exactly three quadratic subfields, and
they close under composition into a Klein four-group:
$\Q(\sqrt{-3})$, $\Q(\sqrt5)$, $\Q(\sqrt{-15})$.
\end{theorem}

That there are exactly three and that they close is the content. Moreover no closed
hyperbolic Dehn filling in the surveyed slope grid carries any of the three: the
arithmetic belongs to the \emph{cusped} object. The manifold's group is a congruence
subgroup, and its continuous spectrum is a single channel with scattering
determinant $\phi(s)=\Lambda_K(s-1)/\Lambda_K(s)$ exactly.

\section{The entrance}\label{sec:entrance}

This is the section the paper rests on: if arriving at an exceptional algebra is
generic, what is \emph{not} generic here?

The manifold's fiber-field side factors through a group of order $360$ at congruence
level $15$. It is tempting---and we record that one of us attempted it---to read
that level as a product of its prime parts and conclude the algebra factorizes.
\textbf{That reading is false.} The congruence-level-$15$ form is a unique
irreducible coupling of the two arithmetic sides: $59$ of $60$ primes falsify the
corresponding $L$-factorization. The two sides interfere; they do not decompose.

\subsection{The two curvature ends}
The cone-manifold family of the figure-eight complement---hyperbolic below cone
angle $2\pi/3$, \emph{Euclidean exactly at} $2\pi/3$, spherical beyond---is due to
Thurston, with the figure-eight cone manifolds worked out by
Hilden--Lozano--Montesinos~\cite{HLM95} and the degenerations by Hodgson and Porti.
\textbf{We claim none of that geometry.} What we add is which arithmetic sits at
each end.

\begin{theorem}\label{thm:ends}
Along that transition the hyperbolic end carries $(\Q(\sqrt{-3}),2T)$ and the
spherical end carries $(\Q(\sqrt5),2I)$, so the dual McKay pair $E_6+E_8$ is
realized as the two stable geometries of a \emph{single} object separated by the
Euclidean wall, rather than as two unrelated faces. At the spherical end the double
cover is the lens space $L(5,2)=S^3/\Z_5$, with $|H_1|=\det(4_1)=5$: the $5$ of that
end is the knot determinant.
\end{theorem}

\begin{scope}\label{sc:asym}
\textbf{The two ends are not symmetric, and the asymmetry runs against the direction
the phrasing invites.} At the hyperbolic end $\pi_1(4_1)\twoheadrightarrow 2T$ is a
genuine group surjection, with exactly two such quotients. At the spherical end
there is \emph{no} surjection at all: $\pi_1(4_1)$ does not surject onto $2I$, nor
onto $A_5$. So the $E_8$ end is \emph{field-level only}, and the two ramified-prime
reductions are symmetric as fields but not as group surjections. (Established by
direct computation; cf.~\cite{Stuebner} for the representation landscape.)
\end{scope}

We state this beside the theorem because a reader who takes ``the two ends carry
$2T$ and $2I$'' to mean two group-level presences has been misled by symmetric
phrasing.

\subsection{Where the uniqueness lives}
Arrival at $E_6$ is generic (\S\ref{sec:suspicious}). The entrance is not.

\begin{theorem}\label{thm:uniq}
Over the infinite metallic family of grammars $R^mL^m$, the golden grammar is the
unique one whose shadow modulus $m^2+4$ reduces to a McKay group. Completeness is a
proof, not a search: writing $|\SL(2,\Z/N)|=N^3\prod_{p\mid N}(1-p^{-2})$ and
bounding $\prod_{p\mid N}(1-p^{-2})\ge\prod_{k=2}^{N}(1-k^{-2})=(N+1)/2N$ gives
$|\SL(2,\Z/N)|\ge N^2(N+1)/2$, already $126>120$ at $N=6$ and increasing. Hence
$|\SL(2,\Z/N)|$ is a McKay-group order only for $N\in\{3,4,5\}$, and $m^2+4\in\{3,4,5\}$
has the unique metallic solution $m=1$.
\end{theorem}

\begin{theorem}[quantization-index law]\label{thm:qil}
The reduction modulo the shadow modulus is an isomorphism exactly when the cusp
conductor is coprime to it, and acquires a kernel line exactly when ramified. For
the golden, $\gcd(4,5)=1$---an isomorphism. For the silver, $\gcd(2,8)=2$---ramified.
\end{theorem}

So the golden's shadow is the McKay group for \emph{two} independent reasons, and
the silver satisfies neither.

\begin{scope}\label{sc:uniqscope}
Three limits, stated with the theorem. \emph{(i)} The uniqueness is proved at the
$E_8$ end; the corresponding statement at the $E_6$ end is not claimed. By
Scope~\ref{sc:asym} the $E_8$-end object is not a quotient of the knot group at all
---which is not a weakness but a statement of what the theorem quantifies over:
grammars and $\SL(2,\Z/N)$, never manifolds. \emph{(ii)} The moduli $\{3,4,5\}$ do
\emph{not} realize the whole exceptional series. $\SL(2,\Z/3)\cong2T$ and
$\SL(2,\Z/5)\cong2I$ give $E_6$ and $E_8$ genuinely, but $E_7=2O$ never occurs:
$|\SL(2,\mathbb{F}_p)|=p(p^2-1)$ is never $48$, and although $|\SL(2,\Z/4)|=48=|2O|$
that group has seven involutions where every finite subgroup of $\SU(2)$ has exactly
one. The order is a coincidence. \emph{(iii)} The filter is finer than the field:
$m=1,4,11$ all give $\Q(\sqrt5)$, but only $m=1$ lands at group level
($120$ versus $5760$ and $1875000$).
\end{scope}

\begin{scope}\label{sc:priorart}
Two literatures bound this claim and we name both. That the golden ratio,
icosahedral symmetry, $2I$ and $E_8$ form a connected cluster is well
established~\cite{Baez,Dechant}, and \S\ref{sec:suspicious} already concedes that
arrival there is not evidence. \textbf{The nearer literature is the other one}: the
metallic grammars $R^mL^m$ are Sturmian morphisms, hence positive automorphisms of
$F_2$, which is Reutenauer's Part~II subject~\cite{Reutenauer}; and matrices attached
to Christoffel words, their group structure and their determinants have been studied
directly~\cite{ChristoffelMatrices}. The claim defended here is neither ``the golden
is connected to $E_8$'' nor a statement about Christoffel matrices as such, but a
\emph{selection} over a family in which each member is evaluated \textbf{at its own
modulus}. We have not found that step in either literature; we record it as
not-found rather than as absent.
\end{scope}

\begin{census}\label{cen:forced}
Of the $43$ links, $26$ are theorems, $6$ exact identities, $5$ no-go results, $1$ a
corollary, $1$ a census---and $4$ are declared choices, at links $3$, $4$, $5$ and
$18$. \textbf{Between link $6$ (the manifold) and link $17$ there is no declared
choice.}
\end{census}

\section{The frame}\label{sec:frame}

Write $\ee_6$ for the Chevalley algebra of type $E_6$ over $\Q$. Inside it sits
$C=\langle x_8,x_{14},x_{16},x_{22}\rangle$, the span of the four binary-tetrahedral
invariants under the principal $\mathfrak{sl}_2$---Klein forms of degrees $8,14,16,22$,
twice the exponents $4,7,8,11$.

\begin{proposition}\label{prop:C}
$C$ is abelian with $\uu(1)^4$ closure, and the four Klein forms are the complete
list of $2T$-invariants in $\ee_6$ under the principal $\mathfrak{sl}_2$.
\end{proposition}

\begin{proposition}\label{prop:mu}
The pencil determinants are exact powers, up to nonzero rational constants, of an
irreducible cubic $\mu$ with $S_3$ Galois group and three real roots. Writing
$K=\Q[\rho]/\mu$, one has the reduced model
\[
   K\cong\Q[x]/(x^3-12x-5),\qquad \disc K=6237=3^4\cdot7\cdot11,
\]
monogenic, totally real, unit rank $2$, and \textbf{class number $h=1$}. The
squarefree part is $7\cdot11=77$, so the quadratic resolvent of the $S_3$ closure is
$\Q(\sqrt{77})$. Only $3,7,11$ ramify; $13,17,19$ are inert; and $2$, $5$ and the
value primes $953$, $1129$, $421493$ are partially split with exactly one degree-one
place each.
\end{proposition}

\begin{remark}
A correction to the record, since the paper would otherwise inherit it. The
resolvent $77$ is a fact about $K$; a companion clause that the golden $5$ ``enters
by ramification'' is not. A ramified prime must divide the field discriminant, and
$5\nmid6237$; in $K$ the prime $5$ is unramified with the same splitting shape as
the value primes, and it is unramified in the whole $S_3$ closure. The $5$ appearing
in $\disc\mu$---like the $13^6$, for $13$ is inert---is borne by the integral model,
not by the arithmetic of the field.
\end{remark}

\section{The cascade}\label{sec:cascade}

\begin{theorem}[first measurement]\label{thm:fmt}
On the plane $\langle x_8,x_{16}\rangle$ the centralizer dimension stratifies as
$12$, $30$, and exactly three lines at $46$; the three lines are the $S_3$-Galois
orbit of the roots of $\mu$. For each such line $K_i$, $z(K_i)\cong\so(10)\oplus\uu(1)$,
while $z(\text{plane})=\so(8)\oplus\uu(1)^2$. Moreover
$\ee_6=(\so(8)\oplus\uu(1)^2)\oplus V_1\oplus V_2\oplus V_3$ with $\dim V_i=16$ and
$[V_i,V_j]=V_k$, and $\ee_6/K_i\cong16\oplus\overline{16}$.
\end{theorem}

The three walls are Galois conjugates of one another---one object seen three ways.

\begin{theorem}[the frame is the magic square]\label{thm:magic}
The algebra built from split octonions and split complexes by the Freudenthal--Tits
construction satisfies the Jacobi identity on all $76{,}076$ unordered basis
triples, has Cartan matrix $E_6$, and admits an explicit isomorphism onto the
Chevalley $\ee_6$ verified on all $3{,}003$ basis pairs with zero discrepancies.
\end{theorem}

This identification is the paper's ceiling: everything derivable at this layer is
standard exceptional-algebra theory~\cite{BartonSudbery,Slansky}.

\begin{theorem}[second measurement]\label{thm:smt}
Over $K$, at the wall point $y^*$, $\dim z(x_1,y^*)=14$ with derived algebra of
dimension $11$ and centre at least $3$, whence
$z(x_1,y^*)=\su(3)\oplus\su(2)\oplus\uu(1)^3$ exactly. The within-$C$ centralizer
ladder is $\{78,46,30,18,14,12\}$: the dimension $26$ is not attained, so the
second measurement skips $\su(5)\oplus\uu(1)$.
\end{theorem}

\section{The real form}\label{sec:realform}

\begin{theorem}[sign-locking]\label{thm:signs}
Every $C$-stabilizing automorphism acts $\pm$-diagonally on the four charges. Two
exact rational trace moments of degree $6$ force $\varepsilon_8\varepsilon_{16}=+1$ and
$\varepsilon_{14}\varepsilon_{22}=+1$, so twelve of the sixteen sign patterns are
impossible; the wall-real pattern is unique, $\varepsilon=(-1,+1,-1,+1)$.
\end{theorem}

The wall is complex in the split frame at all three Galois roots, and the real form
in which the second wall is real is the quaternionic $\ee_{6(2)}$.

A final involution $D_2$---an eleven-flip diagonal on the $27$---carries the ratio
structure: $D_2=\pm\rho_{27}(\sigma_{\chi_-})$ for the \emph{affine} shifted
character, the unique match in a $128$-member family with all $64$ un-shifted
characters failing, and $\{I,D_2,D,D_2D\}$ realizes the wall pair's full $2$-torsion.
From the characterization follow exact $K$-arithmetic invariants, including a norm
law that cubes on the three-dimensional coloured atoms and a sum rule equal to $11$.
\textbf{We record these as abstract invariants of $K$ and as nothing else}; the
primes entering them are not derived further, and the stage at which any comparison
would be made lies outside this paper.

\medskip
\noindent\textbf{Character of \S\S\ref{sec:frame}--\ref{sec:realform}.} Every
statement is a theorem or an exact identity about centralizers, structure constants
or trace moments of an algebra the construction was handed. \textbf{No choice is
made anywhere in them.}

\section{What the construction does not yield}\label{sec:negatives}

We ask the reader to weigh this section equally with the preceding three. A
construction that reaches structure and \emph{cannot} reach magnitude is a different
object from one that has not yet reached magnitude, and only the first is what we
claim.

\subsection{Five sealed predictions, five misses}
Each comparison was pre-registered---prediction, target and decision rule sealed
before computation---and each is reported whatever it returned. The first missed at
$16\sigma$; the second returned the pre-written negative; the third was a
hit-shape but a miss at the deciding tier; the fourth missed in both sectors; and
the fifth, a pre-committed re-run with zero freedom, produced eight misses over
eight pairs. That fifth run met a trigger its own predecessor had sealed in advance
(``an upper-error shrink of about $2\times$ decides sharply'') at $2.21\times$, with
one variant $44$ degrees outside the $1\sigma$ interval.

\subsection{The surface is exhausted, not merely unlucky}
\begin{census}\label{cen:exhausted}
Exactly two target-shaped relations exist in the inventory. Both were fired and both
missed, against the current global fit, in both orderings, with the free phase
profiled over the full circle: one at $3.4\sigma$ on a phase-independent anchor and
therefore unrescuable, the other at $4.7\sigma$.
\end{census}
There is no third relation to try. That is the content of the census.

\subsection{Scale-blindness is a theorem}
\begin{theorem}\label{thm:scale}
The boundary action reduces to $S=-\mathrm{CS}\cdot k-\Vol\cdot\sigma$, and the
object's amphichirality forces $\mathrm{CS}\equiv0$. Hence
$\partial S/\partial k=-\mathrm{CS}\equiv0$ identically: the object is blind to the
level $k$, and $k$-blindness \emph{is} amphichirality.
\end{theorem}
The vanishing deletes the only term that could carry a scale. The silence about
magnitudes is not a gap awaiting a better calculation.

\begin{nogo}\label{ng:torus}
The cusp-torus route is obstructed on two independent legs: amphichirality deletes
the quantized sector of the boundary action, so rational level-one data has no level
to attach to; and the surviving unquantized sector carries no chiral algebra to
decompose, so the quantity one would evaluate is not merely unknown but undefined.
\end{nogo}

\subsection{The declared inputs}
The $18$ algebra steps of Census~\ref{cen:gamma} carry declared input riders:
chirality (three steps), a vacuum direction, and a generation count. On the last,
four clauses travel together: a three \emph{is} object-forced, as
$h^1(M;\mathbf{27})=3$; that three is \emph{not} a generation triple, since no
$\Z/3$ acts on the object and $\mathrm{Isom}(4_1)=D_4$; therefore ``three
generations'' is a \textbf{declared input} wherever a derivation uses it---booked and
priced, but not an output; and the branch structure carrying the ratios is
independent of the preceding two clauses.

\section{Falsifiers}\label{sec:falsifiers}

A construction that cannot fail is not a construction. Four statements would end this
programme, each computable, none requiring a measurement. \textbf{One has already
fired.}

\begin{enumerate}[label=(F\arabic*),leftmargin=2.6em]
\item \emph{Chirality is impossible from any closing.} Open in the falsifying
  direction: chirality is currently an external input, not a proven impossibility.
\item \emph{Every dimensionless invariant the object produces is base-rate
  reachable.} \textbf{Fired.} Of $23$ enumerated invariants, $22$ were graded
  unreachable or base-rate dead; the $23$rd fell three independent ways---a
  parametrization tautology for every phase, an apparent sharpness refuted
  ($0.89\sigma$, not the $7\sigma$ originally claimed), and zero gated hits over the
  object's whole real algebraic tower.
\item \emph{The value layer is frame-relative in every frame.} Known for one pair; a
  proof for all frames closes the value question permanently.
\item \emph{The rank obstruction has no escape.} Proved for the centralizer class;
  the one escape is external, and its canonicity is structurally unreachable by
  orbit invariants.
\end{enumerate}

\begin{scope}
(F2) fires \textbf{for the enumerated space, not in principle}: the atlas is $23$
enumerated invariants, not a proof that none exists. The consequence we accept is
that the construction has no live candidate for a derived dimensionless value.
\textbf{Theorem~\ref{thm:uniq} yields integers, not measured ratios}, and is
therefore not a candidate under (F2) and must not be presented as one.
\end{scope}

\noindent\textbf{What would not falsify it.} A failed comparison with a measured
quantity does not falsify \S\S\ref{sec:genesis}--\ref{sec:realform}; none of those
theorems asserts a measured value. Conversely---and this is the direction that
matters---a \emph{successful} comparison would not confirm them either: arrival at
this structure is generic and the value layer is base-rate dead, so an agreement
would have to survive the same base-rate test everything else failed.
\textbf{The construction is falsifiable as mathematics and, at present, not testable
as physics.}

\appendix

\section{Verification}\label{app:verify}

The empirical failure mode of this programme has been \emph{retrieval}, not
computation: errors have been navigational, while computations reproduced exactly.
Accordingly the paper carries its verification inline, and a reader should never
need an external corpus. Each script accompanying this paper imports nothing
project-internal, uses exact arithmetic in every verdict-bearing comparison, prints
its expected values, and exits non-zero on drift; a single runner executes all of
them and fails if the suite is empty. Proposition~\ref{prop:mu} in particular is
checked end to end, including $h=1$ by Minkowski reduction with an explicit
generator exhibited for \emph{every} prime ideal of norm below the bound---both
primes above $7$ and both above $11$ included, where a single norm-$p$ element would
establish only that a class squares to the identity.

\begin{thebibliography}{9}

\bibitem{Baez} J.~C. Baez, \emph{From the icosahedron to $E_8$},
arXiv:1712.06436.

\bibitem{BartonSudbery} C.~H. Barton and A. Sudbery, \emph{Magic squares and matrix
models of Lie algebras}, Adv. Math. \textbf{180} (2003), no.~2, 596--647.

\bibitem{Dechant} P.-P. Dechant, \emph{The birth of $E_8$ out of the spinors of the
icosahedron}, Proc. R. Soc. A \textbf{472} (2016), no.~2185, 20150504.

\bibitem{HLM95} H.~M. Hilden, M.~T. Lozano and J.~M. Montesinos, \emph{On a
remarkable polyhedron geometrizing the figure eight knot cone manifolds},
J. Math. Sci. Univ. Tokyo \textbf{2} (1995), no.~3, 501--561.

\bibitem{ChristoffelMatrices} \emph{Christoffel matrices and Sturmian
determinants}, arXiv:2409.09824.

\bibitem{MorseHedlund} M. Morse and G.~A. Hedlund, \emph{Symbolic dynamics II.
Sturmian trajectories}, Amer. J. Math. \textbf{62} (1940), 1--42.

\bibitem{Reutenauer} C. Reutenauer, \emph{From Christoffel Words to Markoff
Numbers}, Oxford University Press, 2019.

\bibitem{Slansky} R. Slansky, \emph{Group theory for unified model building},
Phys. Rep. \textbf{79} (1981), 1--128.

\bibitem{Stuebner} M. Stuebner, \emph{Hyperbolic integral homology spheres and
binary icosahedral representations}, arXiv:2502.06488.

\end{thebibliography}

\end{document}
